% ============================================================
%  ch01 —— §0 引言（原文 p.1–p.6）
% ============================================================
\routenote{p.1--p.6}

\begin{center}
{\large\bfseries 摘\quad 要}
\end{center}

Deligne 在 $\PP^1-\{0,\infty\}$ 上构造了一个引人注目的局部系统，它附着于一族 Kloosterman 和。Katz 计算了它的单值性，并问：是否存在一般约化群的 Kloosterman 层？以及在 Langlands 对应下，这些局部系统应当附着于哪些自守形式？受 Gross 与 Frenkel--Gross 工作的启发，我们在几何 Langlands 纲领的框架内找到了一族显式的这类自守形式，甚至一族简单的自守层。我们用这些自守层以统一的方式对任意约化群构造 $\ell$-进 Kloosterman 层，并描述这些 Kloosterman 层的局部与整体单值性。特别地，它们给出具有例外单值群 $G_2$、$F_4$、$E_7$ 与 $E_8$ 的母题 Galois 表示。这也给出任意约化群带野分歧的几何 Langlands 对应的一个例子。

\setcounter{section}{-1}
\section{引言}

\subsection{Kloosterman 和与 Kloosterman 层的回顾}

设 $n$ 是正整数，$p$ 是素数。对每个 $\FF_q\supset\FF_p$ 与 $a\in\FF_q^\times$，$n$ 变量的 \emph{Kloosterman 和}定义为指数和
\[
\Kl_n(a;q):=(-1)^{n-1}\sum_{x_1x_2\cdots x_n=a,\ x_i\in\FF_q}
\exp\Big(\frac{2\pi i}{p}\mathrm{Tr}_{\FF_q/\FF_p}(x_1+x_2+\cdots+x_n)\Big).
\]
Kloosterman 和出现在模形式的 Fourier 系数中。它们满足 Weil 界
\begin{equation}
|\Kl_n(a;q)|\le n q^{(n-1)/2}.
\end{equation}
当 $n=2$ 时，对任意 $a\in\FF_p^\times$，定义角 $\theta(a)\in[0,\pi]$ 使
\[
2p^{\deg(a)/2}\cos\theta(a)=\Kl_2(a;p^{\deg(a)}),
\]
其中 $\deg(a)$ 是 $a$ 在 $\FF_p$ 上的次数。于是当 $\deg(a)$ 趋于 $\infty$ 时，角 $\{\theta(a)\mid a\in\FF_p^\times\}$ 按 Sato--Tate 测度 $\frac{2}{\pi}\sin^2\theta\,d\theta$ 在 $[0,\pi]$ 上等分布。一般地，关于 $\Kl_n(a;q)$ 的类似等分布定理由 Katz [25] 借助 Deligne 在 [11] 中的结果证明。

上述 Kloosterman 和的性质是借助一个层论的化身证明的。让我们回顾这个构造。在 [10] 中，Deligne 考虑图表
\[
\begin{tikzcd}
& \GG_m^n \arrow[dl, "\sigma"'] \arrow[dr, "\pi"] & \\
\GG_a & & \GG_m
\end{tikzcd}
\]
这里 $\GG_m$ 是乘法群，$\GG_a\cong\Aff^1$ 是加法群，而 $\sigma$（相应地 $\pi$）是取 $\GG_m^n$ 的 $n$ 个坐标之和（相应地之积）的映射。设 $\psi:\FF_p\to\mathbb{Q}_\ell(\mu_p)^\times$ 是非平凡特征（这里 $\mu_p$ 是 $p$ 次单位根之集），$\mathrm{AS}_\psi$ 是 $\FF_p$ 上加法群 $\GG_a$ 上相应的 Artin--Schreier 局部系统。Deligne 于是把 Kloosterman 层定义为 $\FF_p$ 上 $\GG_m$ 上如下带 $\mathbb{Q}_\ell(\mu_p)$-系数的层复形：
\[
\Kl_n:=R\pi_!\sigma^*\mathrm{AS}_\psi[n-1].
\]
固定一个嵌入 $\iota:\mathbb{Q}_\ell(\mu_p)\hookrightarrow\CC$，使 $\iota\psi(x)=\exp(2\pi i x/p)$ 对 $x\in\FF_p$ 成立。对任意 $a\in\GG_m(\FF_q)=\FF_q^\times$，记 $\mathrm{Frob}_a$ 为 $a$ 处的几何 Frobenius，它作用在几何茎 $(\Kl_n)_a$ 上。则由 Grothendieck--Lefschetz 迹公式有
\[
\Kl_n(a;q)=\iota\,\mathrm{Tr}\big(\mathrm{Frob}_a,(\Kl_n)_a\big).
\]
在这个意义上，$\Kl_n$ 是 Kloosterman 和族 $\{\Kl_n(a;q)\}_{a\in\FF_q^\times}$ 的层论化身。在 [10, Théorème 7.4, 7.8] 中 Deligne 证明了：

\begin{enumerate}[label=(\arabic*), leftmargin=1.6em, itemsep=.15em]
\item $\Kl_n$ 集中在次数 $0$，且是秩 $n$ 的局部系统。
\item $\Kl_n$ 在 $\{0\}$ 附近驯顺分歧，且单值是只有一个 Jordan 块的幂幺变换。
\item $\Kl_n$ 在 $\{\infty\}$ 附近完全野分歧（即 $\infty$ 处的野惯性群在 $\Kl_n$ 的茎上没有非零不动向量），且 Swan 导子 $\mathrm{Swan}_\infty(\Kl_n)=1$。
\item $\Kl_n$ 是纯权 $n-1$ 的（这蕴涵上文给出的 Weil 界估计）。
\end{enumerate}

在 [25, §8.7] 中，Katz 证明了 Kloosterman 层的唯一性：若 $\GG_m$ 上秩 $n$ 的局部系统 $L$ 满足上面列出的性质 (2) 与 (3)，则 $L$ 在 $\GG_m\otimes_{\FF_p}\bar{\FF}_p$ 上同构于 $\Kl_n$ 经 $\GG_m$ 上平移作用后的结果。此外在 [25, §11] 中，Katz 进一步研究了 $\Kl_n$ 的整体几何单值性。固定 $\GG_m$ 的一个几何点 $\eta$，记
\[
\phi:\pi_1(\GG_m,\eta)\to GL_n(\mathbb{Q}_\ell(\mu_p))
\]
为与 Kloosterman 层 $\Kl_n$ 相伴的单值表示。设 $\phi^{\mathrm{geo}}$ 是 $\phi$ 到 $\pi_1(\GG_m\otimes_{\FF_q}\bar{\FF}_q,\eta)$ 的限制。Katz 定出了 $\phi^{\mathrm{geo}}$ 的像的 Zariski 闭包 $\ckG^{\mathrm{geo}}$：
\begin{equation}
\ckG^{\mathrm{geo}}=
\begin{cases}
\mathrm{Sp}_n, & n \text{ 偶},\\
\mathrm{SL}_n, & n \text{ 奇},\ p \text{ 奇},\\
\mathrm{SO}_n, & n \text{ 奇},\ n\ne 3,\ p=2,\\
G_2, & n=7,\ p=2.
\end{cases}
\end{equation}
借助 Deligne 在 [11] 中的定理，他由此推出 Frobenius 元 $\mathrm{Frob}_a$ 的共轭类按一个「Sato--Tate」测度等分布。

\subsection{本文的动机与目标}

鉴于例外群 $G_2$ 作为整体几何单值群的神秘出现，Katz 问道（[26] p. I-5）：是否所有半单群都会作为 $\GG_m$ 上局部系统的几何单值群出现？换言之，是否存在指数和，其等分布律由任意的单群、特别是例外群支配？

在本文中我们找到这类局部系统的一个统一构造。对任意分裂约化群 $\ckG$，我们将在 $\GG_m=\PP^1\setminus\{0,\infty\}$ 上构造一个 $\ckG$-局部系统 $\Kl_{\ckG}$，其局部分歧与 $\Kl_n$ 类似；而当 $\ckG=GL_n$ 时，我们恢复 Deligne 的 Kloosterman 层 $\Kl_n$。我们将定出其整体几何单值性的 Zariski 闭包（它结果是「大的」），证明该层的纯度并推出等分布律。最后，我们给出关于这类 Galois 表示（或局部系统）唯一性的一个猜想。

为引言的目的，让我们限制到如下情形。设 $G$ 是 $k=\FF_q$ 上分裂、几乎单且单连通的群，或 $G$ 是 $k$ 上的 $GL_n$。设 $\ckG$ 是其在 $\mathbb{Q}_\ell(\mu_p)$ 上的 Langlands 对偶群，即一个分裂、单且伴随的群，或 $GL_n$。

我们构造的动机来自有理函数域 $K=k(t)$ 的 Langlands 对应。Katz [25] 在其研究 Kloosterman 层时已经提出：

\begin{quote}
「……把这些结果与 Langlands 哲学所给出的、用自守形式对这些层所作的猜想性描述相比较，将会是很有意思的。」
\end{quote}

在一系列工作 [20]、[21]、[22]、[15] 中，Gross（部分与 Reeder、与 Frenkel 合作）提出了 $G(\mathbb{A}_K)$ 的一个候选自守表示 $\pi$；在 $G=GL_n$ 的情形，它应当给出 Kloosterman 层 $\Kl_n$。让我们简要回顾他们的工作。

在 [20] 中，Gross 与 Reeder 给出了 $G(k((s)))$ 的表示 $V_\varphi$ 的如下构造。固定 $G$ 的一个 Borel 子群 $B$，记 $U$ 为 $B$ 的幂幺根。记 $I(0):=\{g\in G(k[[s]])\mid g\bmod s\in B\}$ 为 $G(k[[s]])$ 的 Iwahori 子群，$I(1):=\{g\in G(k[[s]])\mid g\bmod s\in U\}$ 为 $I(0)$ 的幂幺根。对任意仿射一般特征 $\varphi:I(1)\to\mathbb{Q}_\ell(\mu_p)$（见 §1.3），令 $V_\varphi:=\mathrm{c\text{-}Ind}_{I(1)\times Z(G)}^{G(k((s)))}(\varphi\otimes 1)$ 为 $G(k((s)))$ 的紧诱导表示。Gross 与 Reeder 证明这个表示不可约且超尖。

对任意整体域 $F$，Gross [21] 设法用迹公式得到了 $G(\mathbb{A}_F)$ 的自守表示 $\pi$ 的重数表达式，这些 $\pi$ 的分歧局部成分或是 Steinberg 表示，或是表示 $V_\varphi$。特别地，当整体域是 $K=k(t)$ 时，这个公式蕴涵：对任意半单单连通群 $G$，存在唯一的尖自守表示 $\pi=\pi(\varphi)$，满足

\begin{enumerate}[label=(\arabic*), leftmargin=1.6em, itemsep=.15em]
\item $\pi$ 在 $\{0,\infty\}$ 之外非分歧；
\item $\pi_0$ 是 $G(k((t)))$ 的 Steinberg 表示；
\item $\pi_\infty$ 是 $G(k((s)))$ 的简单超尖表示 $V_\varphi$，其中 $s=t^{-1}$。
\end{enumerate}

受 Langlands 哲学启发，Gross [22] 提出了如下猜想：自守表示 $\pi=\pi(\varphi)$ 应当对应于 $\GG_m=\PP^1\setminus\{0,\infty\}$ 上的一个 $\ckG$-局部系统 $\Kl_{\ckG}(\varphi)$——我们称之为与 $\ckG$ 及 $\varphi$ 相伴的 Kloosterman 层——它具有与 $\pi$ 平行的如下性质：

\begin{enumerate}[label=(\arabic*), leftmargin=1.6em, itemsep=.15em]
\item $\Kl_{\ckG}(\varphi)$ 是 $\GG_m=\PP^1\setminus\{0,\infty\}$ 上的 $\ckG$-局部系统。
\item $\Kl_{\ckG}(\varphi)$ 在 $\{0\}$ 附近驯顺分歧，且单值是 $\ckG$ 中的正则幂幺元。
\item 与 $\ckG$ 的伴随表示相伴的局部系统 $\Kl^{\mathrm{Ad}}_{\ckG}(\varphi)$（它是 $\PP^1\setminus\{0,\infty\}$ 上秩为 $\dim\ckG$ 的局部系统）在 $\{\infty\}$ 附近完全分歧（即 $\infty$ 处的惯性群没有非零不动向量），且 $\mathrm{Swan}_\infty(\Kl^{\mathrm{Ad}}_{\ckG}(\varphi))=r(\ckG)$，即 $\ckG$ 的秩。
\item 对 $\ckG$ 的任意不可约表示 $V$，相应的局部系统 $\Kl^V_{\ckG}(\varphi)$ 是纯的。
\end{enumerate}

对 $\ckG=GL_n$，层 $\Kl_n$ 具有上述性质。对 $A$、$B$、$C$ 与 $G_2$ 型的 $\ckG$，具有上述性质的 $\ckG$-局部系统此前已由 Katz [26] 用逐情形构造、并以 $\Kl_n$ 为基本构件造出。

在 [15] 中，Gross 与 Frenkel 在复数上构造了这类局部系统的类比物。确切地说，他们定义了 $\PP^1_\CC\setminus\{0,\infty\}$ 上的 $\ckG_\CC$-联络 $\nabla^{\ckG}$（依赖于在 $\check{\mathfrak g}$ 的仿射单根空间中选取非零向量的选择）。这个联络在 $\{0\}$ 处有正则奇点、在 $\infty$ 处有不正则奇点，与上面的性质 (2) 和 (3) 平行。他们进一步证明 $\nabla^{\ckG}(\check X)$ 在上同调意义下刚性 [15, Theorem 1]，并计算了这些联络的微分 Galois 群。对 $A$、$B$、$C$ 与 $G_2$ 型的 $\ckG$，他们验证 $\nabla^{\ckG}$ 与 Katz [26] 构造的联络一致——后者是上述 Katz 的 $\ell$-进 $\ckG$-局部系统的类比物。所有这些都给出强有力的证据，表明 $\nabla^{\ckG}(\check X)$ 应当是猜想中的局部系统 $\Kl_{\ckG}(\varphi)$ 的正确 de Rham 类比物。

[15] 中关于猜想性局部系统 $\Kl_{\ckG}(\varphi)$ 的预测，是我们工作的指导方针。

\subsection{构造方法}

我们对 Kloosterman 层的构造可以概括如下。

我们从上面提到的自守表示 $\pi$ 出发。关键的观察是：$\pi$ 含有一个 Hecke 特征函数 $f_\varphi$，它可以利用双 coset
\[
G(K)\backslash G(\mathbb{A}_K)\Big/I^0\times I^\infty(2)\times\prod_{x\ne0,\infty}G(\mathcal O_x)
\]
的组合学显式写出（§2.1）。这个双 coset 的点是 $\PP^1$ 上带特定水平结构的 $G$-丛的模叠 $\Bun_G(0,2)$ 的有理点，于是我们可以把函数 $f_\varphi$ 升级为此叠上的层 $A_\varphi$。然后我们证明这个层确实是 Hecke 特征层。特别地，我们证明：几何 Hecke 算子作用在 $A_\varphi$ 上所得的特征值定义了 $\PP^1\setminus\{0,\infty\}$ 上的一个 $\ckG$-局部系统——即我们的 Kloosterman 层 $\Kl_{\ckG}(\varphi)$。在技术上，这个结果的证明依赖于如下事实：层 $A_\varphi$ 是某个具有仿射支撑的局部系统的 clean 扩张。

在 $G=GL_n$ 的情形，我们显式算出局部系统 $\Kl_{\ckG}(\varphi)$。对特定的 $\varphi$，与 $\Kl_{\ckG}(\varphi)$ 及 $\ckG=GL_n$ 的标准表示相伴的秩 $n$ 局部系统，结果正是 Deligne 的 Kloosterman 层 $\Kl_n$。这里 Deligne 的图表将作为几何 Hecke 变换的一部分重新出现。

当把 $G$ 换成一类特定的、$\PP^1\setminus\{0,\infty\}$ 上的拟分裂群概形 $\mathbb{G}$ 时，我们的构造同样有效；此外还有一个依赖于额外 Kummer 特征的变体。这些推广看起来正是比较不同群的 Kloosterman 层时所需的自然框架。由于这些推广并不需要额外的论证，我们将在更一般的框架下给出构造。

\subsection{Kloosterman 层的性质}

本文最长的一部分（§4–§6）随后用于研究单伴随群 $\ckG$ 的 Kloosterman 层 $\Kl_{\ckG}(\varphi)$ 的局部与整体性质。

首先，借助 [11]，所期望的纯度性质 (4) 是我们构造的推论。当 $\ckG$ 是伴随群时，我们还可以把 $\Kl^V_{\ckG}(\varphi)$ 归一化为纯权 $0$。于是，像前面一样固定嵌入 $\iota:\mathbb{Q}_\ell(\mu_p)\hookrightarrow\CC$，指数和
\[
\Kl^V_{\ckG}(\varphi;a;q)=\iota\,\mathrm{Tr}\big(\mathrm{Frob}_a,(\Kl^V_{\ckG}(\varphi))_a\big),\qquad a\in\FF_q^\times,\ V\in\Rep(\ckG)
\]
满足 Weil 界
\[
|\Kl^V_{\ckG}(\varphi;a;q)|\le\dim V.
\]
关于 $0$ 处单值性的性质 (2)（见定理 1(2)）将在 §4.3 中证明。关于 $\infty$ 处单值性的性质 (3)（见定理 2）将在 §5 中证明。Swan 导子的计算涉及对仿射 Grassmannian 中某些 Schubert 簇的几何作细致研究，这占据了 §5 的很大一部分。利用 Gross--Reeder 的一个结果 [20]，可以给出 $\infty$ 处单值性的显式描述（推论 2.13）。这些结果合起来表明，局部系统 $\Kl_{\ckG}(\varphi)$ 满足 Gross 所期望的性质 (1)--(4)。

再设
\[
\phi:\pi_1(\PP^1\setminus\{0,\infty\},\eta)\to\ckG(\mathbb{Q}_\ell(\mu_p))
\]
为与 Kloosterman 层 $\Kl_{\ckG}(\varphi)$ 相伴的单值表示，$\phi^{\mathrm{geo}}$ 为其到 $\pi_1(\PP^1\setminus\{0,\infty\}\otimes_k\bar k,\eta)$ 的限制。我们发现（定理 3）：当 $p>2$（且 $\ckG=B_3$ 时 $p>3$）时，$\phi^{\mathrm{geo}}$ 的像的 Zariski 闭包 $\ckG^{\mathrm{geo}}$ 与 Frenkel--Gross 在 [15, Cor. 9, 10] 中算出的 $\nabla^{\ckG}(\check X)$ 的微分 Galois 群 $\ckG^{\nabla}$ 一致；我们在表 1 中列出它。

\begin{table}[htbp]
\centering
\small
\caption{$\Kl_{\ckG}(\varphi)$ 的整体几何单值群}\label{tab:1}
\begin{tabular}{@{}ll@{}}
\toprule
$\ckG$ & $\ckG^{\mathrm{geo}}$\\
\midrule
$A_{2n}$ & $A_{2n}$\\
$A_{2n-1}$, $C_n$ & $C_n$\\
$B_n$, $D_{n+1}$ ($n\ge4$) & $B_n$\\
$E_7$ & $E_7$\\
$E_8$ & $E_8$\\
$E_6$, $F_4$ & $F_4$\\
$B_3$, $D_4$, $G_2$ & $G_2$.\\
\bottomrule
\end{tabular}
\end{table}

利用 Katz 的论证（[25, §3, §13]），这蕴涵关于共轭类 $\phi(\mathrm{Frob}_a)$ 的如下等分布律。

\medskip
\noindent\textbf{推论.} 设 $\ckG$ 单且伴随（此时 $\phi$ 的整个像都落在 $\ckG^{\mathrm{geo}}$ 中）。设 $\ckG^{\mathrm{geo},c}\subset\ckG^{\mathrm{geo}}(\CC)$ 是一个紧实形式（利用嵌入 $\iota:\mathbb{Q}_\ell(\mu_p)\hookrightarrow\CC$ 使 $\ckG^{\mathrm{geo}}(\CC)$ 有意义）。记 $\ckG^{\natural}_{\mathrm{geo},c}$ 为 $\ckG^{\mathrm{geo},c}$ 的共轭类之集。

对 $a\in k^\times$，$\phi(\mathrm{Frob}_a)$ 在 $\ckG^{\mathrm{geo}}(\CC)$ 中的共轭类事实上属于 $\ckG^{\natural}_{\mathrm{geo},c}$。当 $\deg(a)$（即域扩张 $k(a)/k$ 的次数）趋于 $\infty$ 时，共轭类
\[
\{\phi(\mathrm{Frob}_a)\mid a\in k^\times\}\subset\ckG^{\natural}_{\mathrm{geo},c}
\]
按 Haar 测度从 $\ckG^{\mathrm{geo},c}$ 到 $\ckG^{\natural}_{\mathrm{geo},c}$ 的推前而等分布。
\medskip

\subsection{未决问题}

与关于 $\Kl_n$ 的已知结果相比，对 $\Kl_{\ckG}(\varphi)$ 唯一缺少的一环是唯一性。我们证明 $\Kl_{\ckG}(\varphi)$ 在上同调意义下刚性，即
\[
H^1\big(\PP^1,j_{!*}\Kl^{\mathrm{Ad}}_{\ckG}(\varphi)\big)=0.
\]
这给出 $\Kl_{\ckG}(\varphi)$ 的物理刚性的证据：任何其他满足性质 (2) 与 (3) 的 $\PP^1\setminus\{0,\infty\}$ 上的局部系统 $L$，都应当在 $\PP^1\setminus\{0,\infty\}\otimes_k\bar k$ 上同构于 $\Kl_{\ckG}(\varphi)$。我们将在猜想 7.1 与 7.2 中陈述确切的唯一性猜想。

表 1 还提示：表中同一行出现的不同群 $\ckG$ 的 Kloosterman 层应当本质上是同一个（见猜想 7.3）。这可以看作 Kloosterman 层的函子性陈述。

此外，Kloosterman 层的单值性只是在分裂群的情形被详细研究过。在 §7.2 中，我们陈述关于 $\Kl_{^LG}(\varphi)$ 的局部与整体单值性的预测，其中 $\mathbb{G}$ 是 $\PP^1$ 上拟分裂、单且单连通的群概形，在 $\{0,\infty\}$ 之外的所有位处有良好约化，并由一个驯 Kummer 覆盖 $[N]:\PP^1\setminus\{0,\infty\}\to\PP^1\setminus\{0,\infty\}$（$z\mapsto z^N$）分裂。

\subsection{全文结构}

在 §1 中，我们定义 $\PP^1$ 上各种群概形 $\mathbb{G}(m_0,m_\infty)$，它们编码了 $G$-丛模叠的水平结构——自守层就定义在这些模叠上。我们还给出这些模叠的几何描述。在 §2 中，我们构造自守层 $A_\varphi$ 并陈述本文的主要结果。证明在 §4–§6 中给出。§3 专门处理 $G=GL_n$ 的情形，我们在此恢复 Deligne 的经典 Kloosterman 层。在 §7 中，我们陈述关于 Kloosterman 层刚性的猜想，以及关于拟分裂群的 Kloosterman 层的预期。

本文包含四个附录。在 §8 中，我们在拟分裂群的一般框架下证明模空间 $\Bun_G$ 的结构定理。在 §9 中，我们证明外自同构与几何 Satake 等价之间的相容性。这是描述 Kloosterman 层的整体单值像所需要的，而我们未能找到该结果的文献出处。在 §10 中，我们收集关于 $\ckG$ 的拟极小表示的一些事实。在 §11 中，我们分析 $G_2$ 的伴随 Schubert 簇的几何——它不适合其他类型群的统一处理。

\medskip
\noindent\textbf{致谢.} 作者们 2009--2010 年在高等研究院（IAS）逗留期间相遇并开始了本项目。他们感谢 IAS 提供的一流工作条件。
